\chapter{ws-Homeomorphism in Topological Space}
\graphicspath{{Chapter5/Chapter5Figs/EPS/}{Chapter5/Chapter5Figs/}}

\section{Introduction}\label{sec5.1}

The concept of generalized homeomorphism was derived and studied in the year 1991 by Maki et all \cite{Maki5}. N. Nagaveni \cite{Nagaveni5} introduced and investigated rwg homeomorphism in topological space. In the year 2002, M Sheik John \cite{Sheik1} introduced and investigated w-homeomorphism in topological space. The intention of the paper is to analyze and study ws-homeomorphism, ws*-homeomorphism and their relation with some existing homeomorphisms in topological space. Also some of their properties have been verified.

In second section of this chapter, we to introduce and study ws-homeomorphism, ws*-homeomorphism and their relation with some existing homeomorphisms in topological space. Also some of their properties have been investigated

\section{ws-Homeomorphism in Topological Space}\label{sec5.2}

\begin{dfn}\label{defi5.2.1}
A bijective map $\TSh: (\TSP,\tau)\to (\TSQ,\sigma)$ is termed as ws-homeomor\-phism if $\TSh$ is both ws-continuous and ws-open map.
\end{dfn}

\begin{exm}\label{exam5.2.2} 
Take up $\TSP=\TSQ=\{1,2,3\}$, $\tau =\{\TSP, \phi,\{1\},\{2\},\{1,2\}\}$ be a topology on $\TSP$ and $\sigma =\{\TSQ,\phi,\{1\},\{2\},\{1,2\},\{1,3\}\}$ be a topology on $\TSQ$ and $\ws\TSC(\TSP)= \{\phi, \TSP, \{1\}, \{2\},\{3\}, \{1, 3\},\{2, 3\}\}$ And $\ws\TSO(\TSQ) =\{\TSQ, \phi,\{1\},\{2\},\{1,2\},\{1,3\}\}$. Take up $\TSh: \TSP\to \TSQ$ defined by identity map then $\TSh$ is ws-homeomorphism 
\end{exm}

\begin{thm}\label{thm5.2.3}
Each homeomorphism is a ws-homeomorphism but inverse is untrue.
\end{thm}

\begin{proof}
Take up $\TSh: \TSP\to \TSQ$ be a homeomorphism. Then $\TSh$ is bijective, continuous and open map. Since each continuous map is ws-continuous and each open map is ws-open, $\TSh$ is ws-homeomorphism.
\end{proof}

\begin{exm}\label{exam5.2.4}
Take up $\TSP=\TSQ= \{1, 2, 3\}$. Let $\tau= \{\phi, \TSP, \{1\}, \{2\},\break \{1,2\}\}$ be a topology on $\TSP$  and $\sigma=\{\phi, \TSQ, \{1\}, \{2\}, \{1, 2\}, \{1 3\}\}$ be a topology on $\TSQ$ and $\ws\TSC(\TSP)= \{\phi, \TSP, \{1\}, \{2\},\{3\}, \{1,3\},\{2, 3\}\}$. Let $\TSh: (\TSP, \tau )\to (\TSQ, \sigma)$ be a function defined by identity map is ws-homeomorphism but not a homeomorphism as the closed set $\TSH=\{2\}$ in $\TSQ, h^{-1} (\TSH)=\{2\}$ is not a closed in $\TSP$ hence not a continuous function therefore $\TSh$ is not a homeomorphism.
\end{exm}

\begin{thm}\label{thm5.2.5}
If a map $\TSh: \TSP\to \TSQ$ is homeomorphism then the following holds.
\begin{enumerate}[(i)]
\item Each semi-homeomorphism is ws-homeomorphism but inverse is untrue.
\item Each $\alpha$-homeomorphism is ws-homeomorphism but inverse is untrue.
\item Each g\#-homeomorphism is ws-homeomorphism but inverse is untrue.
\item Each *g$\alpha$-homeomorphism is ws-homeomorphism but inverse is untrue.
\item Each g$\xi$*-homeomorphism is ws-homeomorphism but inverse is untrue.
\item Each $\alpha$gp-homeomorphism is ws-homeomorphism but inverse is untrue.
\item Each $\ddot{g}$-homeomorphism is ws-homeomorphism but inverse is untrue.
\item Each regular-homeomorphism is ws-homeomorphism but inverse is untrue
\item Each rb-homeomorphism is ws-homeomorphism but inverse is untrue.
\end{enumerate}
\end{thm}

\begin{proof}
Take up $\TSh: \TSP\to \TSQ$ be a homeomorphism. Then $\TSh$ is bijective, continuous and open map. Since each semi (respectively, $\alpha$, g$^{\text{\#}}$, *g$\alpha$, g$\xi$*, $\alpha$gp, $\ddot{g}$, regular and $rb$) continuous map is ws-continuous and each semi-(respectively, $\alpha$, g, *g$\alpha$, g$\xi$*, $\alpha$gp $\ddot{g}$, regular and $rb$) -open map is ws-open, $\TSh$ is ws-homeomorphism.
\end{proof}

\begin{exm}\label{exam5.2.6}
In Example \ref{exam5.2.4}, $\TSh$ is ws-homeomorphism but not a semi (respectively, $\alpha$, g$^{\#}$, *g$\alpha$, g$\xi$*, $\alpha$gp, $\ddot{g}$, regular and $rb$) homeomorphism as the closed set $\TSH=\{2\}$ in $\TSQ$, $h^{-1} (\TSH)=\{2\}$ is not a semi (respectively, $\alpha$, g\#, *g$\alpha$, g$\xi$*, $\alpha$gp, $\ddot{g}$, regular and $rb$) closed in $\TSP$.
\end{exm}

\begin{thm}\label{thm5.2.7}
Each g\#s - homeomorphism is ws-homeomorphism but inverse is untrue.
\end{thm}

\begin{proof}
Take up $\TSh: (\TSP, \tau) \to (\TSQ, \sigma)$ be a homeomorphism. Then $\TSh$ is bijective, continuous and open map. Since each g\#s - continuous map is ws-continuous and each g\#s- open map is ws-open, $\TSh$ is ws-homeomorphism.
\end{proof}

\begin{exm}\label{exam5.2.8}
Take up $\TSP=\TSQ= \{1, 2, 3, 4\}$. Let $\tau= \{\phi, \TSP, \{1\}, \{2\}, \{1, 2\},\break \{1,2,3\}\}$ be a topology on $\TSP$ and a $\sigma=\{\phi,\TSQ,\{1\},\{1,2\},\{1,2,3\}\}$ be topology on $\TSQ$, and $\ws\TSC(\TSP)=\{\TSP,\phi,\{1\},\{2\},\{3\},\{4\},\{1,3\},\{1,4\},\{2,3\},\{2,4\},\{3,4\},\{1,2,4\},\break \{1,3,4\},\{2,3,4\}\}$, $\text{g\#s}\TSC(\TSP)=\{\TSP, \phi,\{2\},\{3\},\{4\},\{2,3\},\{2,4\},\{3,4\},\{2,3,4\}\}$,\break $\ws\TSO(\TSQ)=\{\TSQ, \phi, \{1\},\{2\},\{3\},\{1,2\}, \{1,3\},\{1,4\},\{2,3\},\{1,2,3\}, \{1,2,4\}, \{1,3,4\}\},$ g\#s$\TSO(\TSQ)=\{\TSQ,\phi,\{1\},\{1,2\},\{3,4\},\{1,4\},\{1,2,3\}\{1,2,4\},\{1,3,4\}\}$ Let $\TSh: (\TSP, \tau )\break \to (\TSQ, \sigma)$ be a function defined by $\TSh(1)= 1$, $\TSh(2)=2$, $\TSh(3)=3$, $\TSh(4)=2$, is ws-homeomorphism but not a homeomorphism hunction as the open set $\TSH=\{3,4\}$ in $\TSQ$, $h^{-1} (\TSH)=\{1,2\}$ is not a g\#s-closed in $\TSP$ hence not a continuous function therefore not a g\#s-homeomorphism.
\end{exm}

\begin{thm}\label{thm5.2.9}
If a map $\TSh: \TSP\to \TSQ$ is homeomorphism then
\begin{itemize}
\item[(iv)] If $\TSh$ is ws-homeomorphism then it is gspr-homeomorphism but inverse is untrue.
\item[(v)] If $\TSh$ is ws-homeomorphism then it is gsp-homeomorphism but inverse is untrue.
\item[(vi)] If $\TSh$ is ws-homeomorphism then it is rgb-homeomorphism but inverse is untrue.
\end{itemize}
\end{thm}

\begin{proof}
Take up $\TSh: (\TSP,\tau)\to (\TSQ, \sigma)$ be ws-homeomorphism. Then $\TSh$ is both ws-continuous and ws-open map. This implies that $\TSh$ is both gspr continuous and gspr open map. Therefore $\TSh$ is gspr-homeomorphism.

Similarly we can prove (ii) and (iii).
\end{proof}

\begin{exm}\label{exam5.2.10}
Take up $\TSP=\TSQ= \{1, 2, 3,4\}$. Let $\tau= \{\phi, \TSP, \{1\}, \{2\}, \{1, 2\},\break \{1,2,3\}\}$ be a topology on $\TSP$ and a $\sigma=\{\phi,\TSQ,\{1,2\},\{3,4\}\}$ be topology on $\TSQ$ and $\ws\TSC(\TSP)=\gsp\TSC(\TSP)=\{\TSP,\phi,\{1\},\{2\},\{3\},\{4\},\{1,3\},\{1,4\},\{2,3\},\{2,4\},\{3,4\},\break \{1,2,4\},\{1,3,4\},\{2,3,4\}$, $\gspr\TSC(\TSP)=\rgb\TSC(\TSP)=\{\TSP(\TSP)\}$, $\ws\TSO(\TSQ)=\{\phi,\TSQ,\{1,2\},\break \{3,4\}\}$, $\gspr\TSC(\TSQ)=\rgb\TSC(\TSQ)= \gsp\TSC(\TSQ)=\{\TSP(\TSQ)\}$, and Take up $\TSh: \TSP\to \TSQ$ defined by $\TSh(1)=1$, $\TSh(2)= 2$, $\TSh(3)=3$, $\TSh(4)=2$ then $\TSh$ is gsp-homeomorphism(respectively, gsp-homeomorphism, rgb-homeomorphism). But not ws-homeomorphism, as open set $\TSF= \{1\}$ in $\TSP$, then, $h^{-1} (\TSF) = \{1\}$ is not ws-open set in $\TSQ$.
\end{exm}

\begin{rem}\label{rem5.2.11}
Following example shows that ws-homeomorphism is independent of g, w, $\beta$, gp, g$\alpha$, swg, rwg, wg, g*p, gw$\alpha$, $\alpha${*}{*}g, swg*, gr, w$\alpha\hat{g}$, g$\alpha$*, rg, rg$\alpha$ R*, rgw, wgr$\alpha$, gprw, pgr$\alpha$, g$\alpha${*}{*}, $\beta$wg{*}{*}, mildly g-homeomorphisms.
\end{rem}

\begin{exm}\label{exam5.2.12}
Take up $\TSP=\TSQ= \{1, 2, 3\}$. Let $\tau= \{\phi, \TSP, \{1\}, \{2\}, \{1, 2\}\}$ be a topology on $\TSP$ and $\sigma=\{\phi, \TSQ, \{1\},\{2,3\}\}$ be a topology on $\TSQ$ and $\ws\TSC(\TSP)=\{\TSP, \phi,\{1\}, \{2\}, \{3\},\{2,3\},\{1, 3\}\}$, Let $\TSh: (\TSP, \tau)\to (\TSQ, \sigma)$ be a function defined by $\TSh(1)=1$, $\TSh(2)= 1$, $\TSh(3)=1$ is ws-homeomorphism but not g, w, $\beta$, gp, g$\alpha$, swg, rwg, wg, g*p, gw$\alpha$, $\alpha${*}{*}g, swg*, gr, w$\alpha\hat{g}$, g$\alpha$*, rg, rg$\alpha$, R*, rgw, wgr$\alpha$, gprw, pgr$\alpha$, g$\alpha${*}{*}, $\beta$wg{*}{*}, mildly g-homeomorphisms as the closed set $\TSH=\{1\}$ in $\TSQ$, $h^{-1}(\TSH)=\{1\}$ is not a g(respectively w, $\beta$, gp, g$\alpha$, swg, rwg, wg, g*p, gw$\alpha$, $\alpha${*}{*}g, swg*, gr, w$\alpha\hat{g}$, g$\alpha$*, rg, rg$\alpha$, R*, rgw, wgr$\alpha$, gprw, pgr$\alpha$, g$\alpha${*}{*}, $\beta$wg{*}{*}, mildly g)closed set in $\TSP$.
\end{exm}

\begin{exm}\label{exm5.2.13}
Take up $\TSP=\TSQ= \{1, 2, 3\}$. Let $\tau= \{\phi, \TSP, \{1\}, \{2,3\}\}$ be a topology on $\TSP$ and $\sigma=\{\phi, \TSQ, \{1\},\{2\}, \{1,2\}\}$ be a topology on $\TSQ$ and $\ws\TSC(\TSP)= \{\phi, \TSP,\{1\},\{2,3\}\}$. Let $\TSh: (\TSP, \tau )\to (\TSQ, \sigma)$ be a function defined by $\TSh(1)=3$, $\TSh(2)= 2$, $\TSh(3)=3$, is g, w, $\beta$, gp, g$\alpha$, swg, rwg, wg, g*p, gw$\alpha$, $\alpha${*}{*}g, swg*, gr, w$\alpha\hat{g}$, g$\alpha$*, rg, rg$\alpha$, R*, rgw, wgr$\alpha$, gprw, pgr$\alpha$, g$\alpha${*}{*}, $\beta$wg{*}{*}, mildly g-homeomorphisms but not ws-homeomorphism as the closed set $\TSH=\{3\}$ in $\TSQ$, $h^{-1} (\TSH)=\{1,3\}$ is not a ws-closed in $\TSP$.
\end{exm}

\begin{thm}\label{thm5.2.14}
Take up $\TSh: \TSP\to \TSQ$ be bijective and ws-continuous map. Henceforth following results are one and the same.
\begin{enumerate}[(1)]
\item $\TSh$ is a ws-open map
\item $\TSh$ is a ws-homeomorphism
\item $\TSh$ is a ws-closed map
\end{enumerate}
\end{thm}

\begin{proof}
(1) $\Rightarrow$ (2)

Take up $\TSh: \TSP\to \TSQ$ be bijective, ws-continuous map and ws-open map. Therefore $\TSh$ is a ws-homeomorphism.

(2) $\Rightarrow$ (3)

Take up $\TSh: \TSP\to \TSQ$ be a ws-homeomorphism. Let $\TSG$ be any closed set $\in \TSP$ and hence $\TSG^{\text{c}}$ is open set in $\TSP$. Then $\TSh(\TSG^{\text{c}})$ is open in $\TSQ$. But $\TSh(\TSG^{\text{c}})=\TSh(\TSG)^{\text{c}}$. Therefore $\TSh$ is a ws-closed map.

(3) $\Rightarrow$ (1)

Take up $\TSh: \TSP\to \TSQ$ be bijective, ws-continuous and ws-closed map. Assume $\TSG$ be any open set $\in \TSP$ and hence $\TSG^{\text{c}}$ is closed $\in \TSP$. Then $\TSh(\TSG^{\text{c}})$ is closed $\in \TSQ$. But $\TSh(\TSG^{\text{c}})=\TSh(G)^{\text{c}}$. Therefore $\TSh$ is a ws-open map.
\end{proof}

\begin{rem}\label{rem5.2.15}
The composition of two ws-homeomorphism is not a ws-homeomor\-phism in general as seen in example \ref{exam5.2.16}.
\end{rem}

\begin{exm}\label{exam5.2.16}
Take up $\TSP=\TSQ=\TSR= \{1, 2, 3\}$. Let $\tau= \{\TSP, \phi, \{1\}, \{2\}, \{1,2\},\break \{1,3\}\}$ be a topology on $\TSP$, $\sigma= \{\TSQ, \phi, \{1\}, \{2\}, \{1,2\}\}$ be a topology on $\TSQ$ and $\eta= \{\TSZ, \phi, \{1\}, \{2,3\}\}$ be a topology on $\TSR$ and $\ws\TSC(\TSP)= \{\TSP, \phi, \{2\}, \{3\}, \{1, 3\}, \{2, 3\}\}$, $\ws\TSC(\TSQ)= \{\TSQ, \phi, \{1\},\{2\}, \{3\}, \{1,3\}, \{2, 3\}\}$ $\ws\TSC(\TSR)= \{\TSR, \phi, \{1\}, \{2, 3\}\}$. Define $\TSh: (\TSP, \tau) \to (\TSQ,\sigma)$ by $\TSh(1)=1$, $\TSh(2) =2$, $\TSh(3)=3$ and $\TSg: (\TSQ,\sigma)\to (R, \eta)$ be a function defined by $\TSh(1)=1$, $\TSh(2)= 1$, $\TSh(3)=1$. Both $\TSh$ and $\TSg$ are ws-homeomorphism but their composition go$\TSh: (\TSP, \tau)\to (\TSR, \eta)$ defined by identity function is not a ws-homeomorphism as the closed set $\TSH=\{1\}$ in $(\TSR, \eta)$, but $(goh)^{-1} (\TSH)=\{1\}$ is not ws-closed in $\TSP$.
\end{exm}

\begin{dfn}\label{defi5.2.17}
A bijective map $\TSh: \TSP\to \TSQ$ is called a ws*-homeomorphism if $\TSh$ and $h^{-1}$ are ws-irresolute. Where ws*-$\TSh(\TSP, \tau)$ denote the set of all ws-homeomorphisms of $\TSP$ on to itself.
\end{dfn}

\begin{thm}\label{thm5.2.18}
Each ws*-homeomorphism is a ws-homeomorphism but inverse is untrue.
\end{thm}

\begin{proof}
Take up $\TSh: \TSP\to \TSQ$ be a ws*-homeomorphism. Then $\TSh$ is bijective, $\TSh$ and $h^{-1}$ are ws-irresolute. Since each ws-irresolute map is ws-continuous then $\TSh$ and $h^{-1}$ are ws-continuous. Also $h^{-1} : \TSQ \to \TSP$ is a ws-continuous map is one and the same. to ws-open map. Thus $\TSh$ is bijective, ws-continuous and ws-open map. Hence $\TSh$ is a ws-homeomorphism.
\end{proof}

\begin{exm}\label{exam5.2.19}
Take up $\TSP=\TSQ= \{1, 2, 3\}$. Let $\tau= \{\phi, \TSP, \{1\}, \{2\}, \{1,2\},\{1,3\}\}$ be a topology on $\TSP$ and $\sigma=\{\phi, \TSQ, \{1\}, \{2\},\{1,2\}\}$ be a topology on $\TSQ$ and $\ws\TSC(\TSP)\break = \{\phi, \TSP, \{1\},\{3\}, \{1,3\},\{2,3\}\}$, $\SCL(\TSQ) =\{\phi, \TSQ, \{1\},\{2\},\{3\},\{2, 3\},\{1,3\}\}$,\break $\ws\TSC(\TSQ)= \{\phi, \TSQ, \{1\},\{2\},\{3\}, \{1,3\},\{2,3\}\}$ Let $\TSh: (\TSP, \tau )\to (\TSQ, \sigma)$ be a function defined by identity function is ws-continuous but not a ws*-continuous function as the semi-closed set $\TSF=\{1\}$ in $\TSQ$, $h^{-1}(\TSF)=\{1\}$ not a ws-closed in $\TSP$.
\end{exm}

\begin{thm}\label{thm5.2.20}
Each ws*-homeomorphism is a gspr-homeomorphism (resp gsp-homeomorphism, rgb-homeomorphism) but inverse is untrue.
\end{thm}

\begin{proof}
Proof is obtained from the previous results, definations and evidence each ws*-homeomorphism is a ws-homeomorphism and each ws-homeomorphism is a gspr-homeomorphism (resp gsp-homeomorphism, rgb-homeomorphism)
\end{proof}

\begin{exm}\label{exam5.2.21}
Take up $\TSP=\TSQ= \{1, 2, 3, 4\}$. Let $\tau= \{\phi, \TSP, \{1\}, \{2\}, \{1, 2\},\break \{1,2,3\}\}$ be a topology on $\TSP$ and $\sigma=\{\phi,\TSQ,\{1,2\},\{3,4\}\}$ be a topology on $\TSQ$ and $\ws\TSC(\TSP)= \gsp\TSC(\TSP)=\{\TSP, \phi,\{1\},\{2\},\{3\},\{4\},\{1,3\},\{1,4\},\{2,3\},\{2,4\},\{3,4\},\break \{1,2,4\},\{1,3,4\},\{2,3,4\}\}$, $\gspr\TSC(\TSP)=\rgb\TSC(\TSP)=\{\TSP(\TSP)\}$, $\ws\TSO(\TSQ)=\{\phi,\TSQ,\{1,2\},\break \{3,4\}\}$, $\gspr\TSC(\TSQ)=\rgb\TSC(\TSQ)=\gsp\TSC(\TSQ)=\{\TSP(\TSQ)\}$. Take $\TSh: \TSP\to \TSQ$ defined by $\TSh(1)=1$, $\TSh(2)= 2$, $\TSh(3)=3$, $\TSh(4)=2$ then $\TSh$ is gsp-homeomorphism(Respectively, gsp-homeomorphism, homeomorphism). But not ws-homeomorphism, as open set $\TSH= \{1\}$ in $\TSP$, then, $h^{-1} (\TSH) = \{1\}$ which is not ws-open set in $\TSQ$.
\end{exm}

\begin{rem}\label{rem5.2.22} 
From the above discussions and known facts, the relation between ws-homeomorphism with some existing homeomorphisms in topological space is shown in
following figure.
\begin{figure}[!ht]
\centering
\includegraphics[scale=.72]{1.eps}
\end{figure}
\end{rem}

\begin{thm}\label{thm5.2.23} 
The composition of two ws*-homeomorphisms is also a ws*-{\break}homeomor\-phism.
\end{thm}

\begin{proof}
Take up $\TSh: \TSP\to \TSQ$ and $\TSg: \TSQ\to \TSR$ be two ws*-homeomorphisms. Now we have to prove goh: $\TSQ\to \TSR$ be ws*-homeomorphism. Let $\TSN$ be any closed set in $\TSR$. Seeing that $\TSg$ is ws-irresolute, $\TSg(\TSN)$ is ws-closed $\in \TSQ$. Also $\TSh$ is ws-irresolute, $h^{-1} (g^{-1} (\TSN))$ is ws-closed in $\TSP$. But $h^{-1} (g^{-1}(\TSN))= (goh)^{-1}(\TSN)$. Therefore goh is ws-irresolute. Take $\TSW$ be any closed set in $\TSP$. Since $\TSh$ is ws-irresolute, $\TSh(\TSW)$ is ws-closed in $\TSQ$ and also $\TSg(\TSh(\TSW))$ is ws-closed in $\TSR$. But (goh)$(\TSW)=\TSg(\TSh(\TSW))$. Therefore (goh) is ws-irresolute. Hence (goh) is ws*-homeomorphism.
\end{proof}

\begin{thm}\label{thm5.2.24}
Under the composition of maps, the set ws*-$\TSh(\TSP, \tau)$ is a group.
\end{thm}

\begin{proof}
Take up * be a binary operation, $*: \ws * \TSh(\TSP, \tau) \to  \ws * \TSh(\TSP, \tau)$ defined by $\TSh * \TSg = \text{goh}$ for all $\TSh$, $\TSg \in \ws * \TSh(\TSP, \tau)$. We have, $\text{goh} \in \ws * \TSh(\TSP,\tau)$. Also the associative property holds in composition of maps. The identity map $\TSI : (\TSP, \tau) \to  (\TSP, \tau)$ which $\in \ws *\TSh(\TSP, \tau)$, acts as the identity element. For the inverse element, if $\TSh \in \ws * \TSh(\TSP,\tau)$ then $h^{-1} \in \ws * \TSh(\TSP,\tau) \ni: hoh^{-1}=h^{-1}oh=\TSI$. Therefore, under the composition of maps, the set $\ws *-\TSh(\TSP,\tau)$ is a group.
\end{proof}

